\section{Open Questions and Risks}
\label{sec:open}

\paragraph{Feature identity across time.} The time axis comes from the trajectory
(Sec.~\ref{sec:temporal}), and the whole construction assumes a feature means the
same thing at step $t$ and step $t{+}1$. If a feature drifts in meaning, every
causal edge built on its time series is suspect. We still owe this an explicit
stability check on GraphCast rollout data.

\paragraph{Model dynamics vs.\ atmospheric dynamics.} Discovery on internal
features recovers the causal structure of the model's computation. Whether that
matches the atmosphere's causal structure is a separate empirical question. The
interpretability framing (``what did the model learn?'') legitimately targets the
model's structure, but every claim should say which of the two it is about.

\paragraph{No ground truth in GraphCast.} SAVAR lets us score recovered graphs;
GraphCast never will. Validation there has to lean on known physical relationships
(the diurnal cycle driving near-surface temperature, say) and on interventional
consistency in the style of~\cite{macmillan2025}: an edge $A \to B$ produced by
discovery should survive an actual intervention on $A$.

\paragraph{Polysemanticity.} In every SAVAR variant tested, the best-aligned
features remain polysemantic, and a polysemantic feature is an ambiguous causal
variable. Possible responses include stricter monosemanticity criteria, per-region
SAEs, or grouping features before discovery. The heterogeneous-dynamics experiment
(Sec.~\ref{sec:homog}) will show whether specialization appears once the task
rewards it; if it does not, the grouping problem moves to the center of the
project.

\paragraph{Orientation of contemporaneous edges.} The aliasing risk itself is now
quantified: Sec.~\ref{sec:subsample} shows plain PCMCI collapsing under
subsampling while PCMCI+ recovers the aliased adjacencies. Time priority still
orients every lagged edge, and a time-unrolled lagged graph is a DAG, so feedback
loops such as X2$\to$X0 stay identifiable. The unsolved part is direction at
$\tau=0$: ParCorr orients only $8$--$30\%$ of the contemporaneous edges it finds,
despite innovations skewed enough to make the directions identifiable in
principle. A LiNGAM-style orientation step on top of the PCMCI+ adjacencies is the
natural fix and is not yet implemented. On GraphCast at 6-hour cadence, sub-step
atmospheric pathways will alias, and the sweep in Table~\ref{tab:subsample} is a
quantitative preview of what that costs.

\paragraph{Weak-signal regimes.} The fine-cadence one-step task leaves the
forecaster barely above climatology, capping the effect size any ablation can
show (Sec.~\ref{sec:vpd}). Conclusions there rest on which components matter and
where, never on how much. GraphCast at 6 hours is far above climatology, so this
is a SAVAR-specific worry for now, but it will return at long lead times as
rollout skill decays.

\paragraph{Overlapping modes.} Every SAVAR variant so far uses spatially disjoint
mode maps, so the pixel-to-mode projection is clean. Overlapping modes (closer
centers, wider spatial profiles) would degrade the projection and dilute the SAE
features. The knob exists in the generator and has not been turned.

\paragraph{Method disagreement on GraphCast.} If PCMCI+, TSCI, and DYNOTEARS
disagree on GraphCast features, some adjudication rule is needed, and it should be
calibrated on SAVAR. Section~\ref{sec:methodcompare} is the start of one: on
activation-derived series, weight PCMCI+ most, and treat edges only DYNOTEARS
reports with caution.
